% ==============================================================================
% T0-Theorie: Dokument 161
% Titel: T0-Theorie und Coherent Ising Machines
% Autor: Johann Pascher
% Datum: Februar 2026
% ==============================================================================

\section*{Abstract}

Wir analysieren die strukturellen Parallelen zwischen der T0 Zeit-Masse-Dualitätstheorie 
und der Coherent Ising Machine (CIM), einem photonischen Analogcomputer zur Lösung 
kombinatorischer NP-harter Optimierungsprobleme. Es wird gezeigt, dass beide Systeme 
auf demselben fundamentalen Prinzip basieren: \emph{Optimierung als physikalische 
Relaxation in ein geometrisches Minimum} -- nicht als algorithmische Iteration.
Die sechs identifizierten strukturellen Parallelen umfassen Determinismus, Grundzustand, 
Qubit-Struktur, Kopplungsmatrix, fraktale Korrekturen und toroidale Topologie.
Die zentrale Konsequenz ist das Konzept einer \emph{T0-Ising Machine}, bei der alle 
Kopplungen $J_{ij}$ nicht extern programmiert werden, sondern aus dem intrinsischen 
Zeitfeld $\Tfield$ geometrisch emergieren.

\clearpage

\section{Einführung}
\label{sec:intro}

\subsection{Motivation}

Die Coherent Ising Machine (CIM) ist ein optischer Analogcomputer, der kombinatorische 
Optimierungsprobleme durch physikalische Relaxation in den Grundzustand eines 
Ising-Hamiltonians löst \cite{inagaki_cim_2016, mcmahon_cim_2016}. Das fundamentale 
Prinzip -- das System ``fällt'' physikalisch in seine Lösung -- steht in direktem 
konzeptuellem Zusammenhang mit der T0-Theorie \cite{pascher_t0_2025}, die alle 
Naturkonstanten aus einem einzigen geometrischen Parameter $\xipar = 4/30000$ ableitet.

Das vorliegende Dokument untersucht diese Verbindung systematisch. Wir zeigen, dass 
T0 nicht nur eine konzeptuelle Parallele zur CIM liefert, sondern eine fundamental 
stärkere Formulierung ermöglicht: eine T0-basierte Ising Machine, bei der die 
Kopplungsmatrix $J_{ij}$ geometrisch aus dem Massenfeld $\mfield$ emergiert -- 
nicht extern programmiert werden muss.

\subsection{Struktur}

Abschnitt~\ref{sec:ising} beschreibt das Ising-Modell als universelle 
Optimierungssprache. Abschnitt~\ref{sec:cim} erklärt das Funktionsprinzip der CIM.
Abschnitt~\ref{sec:t0} fasst die relevante T0-Grundstruktur zusammen.
Abschnitt~\ref{sec:parallels} analysiert die sechs strukturellen Parallelen systematisch.
Abschnitt~\ref{sec:t0machine} formuliert das Konzept der T0-Ising Machine.
Abschnitt~\ref{sec:questions} diskutiert offene Forschungsfragen.

\clearpage

\section{Das Ising-Modell als universelle Optimierungssprache}
\label{sec:ising}

\subsection{Der Hamiltonoperator}

Das Ising-Modell aus der Festkörperphysik beschreibt magnetische Spins 
$\sigma_i \in \{-1, +1\}$, die über einen Hamiltonoperator wechselwirken:

\begin{equation}
    H_{\text{Ising}} = -\frac{1}{2} \sum_{i \neq j} J_{ij} \, \sigma_i \sigma_j 
    - \sum_{i} h_i \, \sigma_i
    \label{eq:ising_hamiltonian}
\end{equation}

Dabei sind $J_{ij}$ die Kopplungsstärken zwischen Spin $i$ und $j$, und $h_i$ 
externe Felder (Bias-Terme). Der Grundzustand -- die Spin-Konfiguration mit minimalem 
$H_{\text{Ising}}$ -- entspricht der optimalen Lösung des codierten Problems.

\subsection{QUBO: Die universelle Abbildung}

Jedes NP-harte kombinatorische Problem lässt sich als Minimierung des Ising-Hamiltonians 
\eqref{eq:ising_hamiltonian} formulieren, via \emph{Quadratic Unconstrained Binary 
Optimization} (QUBO):

\begin{table}[h]
    \centering
    \caption{Beispiele für die Abbildung NP-harter Probleme auf Ising-Hamiltonians}
    \label{tab:qubo_examples}
    \resizebox{\textwidth}{!}{\begin{tabular}{@{}ll@{}}
        \toprule
        Optimierungsproblem & Ising-Entsprechung \\
        \midrule
        Travelling Salesman Problem & Kopplungen $J_{ij}$ als Distanzmatrix \\
        MaxCut (Graphenpartitionierung) & Spins repräsentieren Knotenzuordnungen \\
        Portfolio-Optimierung & Spins als Aktienauswahl (kaufen/nicht kaufen) \\
        Proteinstrukturvorhersage & Spins als Torsionswinkel-Konfigurationen \\
        RSA-Faktorisierung & Spins als Bitdarstellung der Faktoren \\
        \bottomrule
    \end{tabular}}
\end{table}

Diese Abbildung ist keine Approximation, sondern exakt. Wer den Grundzustand von 
\eqref{eq:ising_hamiltonian} findet, hat das Optimierungsproblem gelöst.

\begin{important}[Zentrales Prinzip]
    Kombinatorische Optimierung $\equiv$ Suche nach dem Grundzustand des Ising-Hamiltonians. 
    Dies ist der konzeptuelle Kern sowohl der CIM als auch -- wie wir zeigen werden -- 
    der T0-Theorie.
\end{important}

\clearpage

\section{Coherent Ising Machine: Funktionsprinzip}
\label{sec:cim}

\subsection{Physikalische Implementierung}

Eine Coherent Ising Machine (CIM) implementiert das Ising-Modell physikalisch durch 
ein Netzwerk zeitgemultiplexter degenerierter optischer parametrischer Oszillatoren 
(DOPO) \cite{yamamoto_cim_2017, honjo_100k_cim_2021}:

\begin{itemize}
    \item \textbf{Spin} $\sigma_i = +1$: Laserpuls mit Phase $0°$
    \item \textbf{Spin} $\sigma_i = -1$: Laserpuls mit Phase $180°$
    \item \textbf{Kopplung} $J_{ij}$: Elektronisches Feedback-Netzwerk zwischen Pulsen
    \item \textbf{Dynamik}: Physikalische Relaxation in das Energieminimum
\end{itemize}

Das System entwickelt sich deterministisch gemäß den Feldgleichungen des 
parametrischen Oszillators -- und findet dabei automatisch (oder mit hoher 
Wahrscheinlichkeit) den Grundzustand des Hamiltonians \eqref{eq:ising_hamiltonian}.

\subsection{Skalierung und Stand der Technik}

\begin{table}[h]
    \centering
    \caption{Vergleich: Klassischer Rechner vs.\ Coherent Ising Machine}
    \label{tab:cim_vs_classical}
    \resizebox{\textwidth}{!}{\begin{tabular}{@{}lll@{}}
        \toprule
        Eigenschaft & Klassischer Computer & Coherent Ising Machine \\
        \midrule
        Parallelismus & Sequenziell / begrenzt & Massiv parallel (alle Spins simultan) \\
        Lösungsparadigma & Algorithmus iteriert & Physikalische Relaxation \\
        Skalierung & Exponentiell schwieriger & Natürlich durch Lichtpropagation \\
        Energiebedarf & Hoch bei NP-harten Prob. & Potenziell Größenordnungen niedriger \\
        Spins (Stand 2024) & -- & Bis zu 100.000 optische Spins \\
        \bottomrule
    \end{tabular}}
\end{table}

\subsection{Determinismus hinter scheinbarer Stochastik}
\label{subsec:cim_determinism}

Eine CIM scheint stochastisch zu suchen: verschiedene Läufe liefern manchmal 
verschiedene Lösungen. Aber dies ist eine Oberflächen-Eigenschaft: physikalisch 
folgt das System \emph{deterministisch} den Feldgleichungen des OPO. Die scheinbare 
Zufälligkeit ist ausschließlich Unwissenheit über die genauen Anfangsbedingungen 
des Quantenvakuums.

Dies ist die erste und grundlegendste Parallele zur T0-Theorie.

\clearpage

\section{T0-Grundstruktur: Relevante Konzepte}
\label{sec:t0}

\subsection{Die fundamentale Bindungsbedingung}

Die T0-Theorie postuliert Zeit-Masse-Dualität als fundamentale geometrische Constraint:

\begin{equation}
    T \cdot m = 1 \quad \text{(in natürlichen Einheiten)}
    \label{eq:t0_constraint}
\end{equation}

Das intrinsische Zeitfeld eines Teilchens mit Masse $m$ ist:

\begin{equation}
    \Tfield = \frac{\hbar}{m(x,t) \cdot c^2}
    \label{eq:t0_timefield}
\end{equation}

Zeit und Masse sind keine unabhängigen Größen -- sie sind reziproke Aspekte einer 
einzigen geometrischen Realität.

\subsection{Der geometrische Ursprungsparameter}

Alle Naturkonstanten emergieren aus einem einzigen dimensionslosen geometrischen 
Parameter \cite{pascher_xi_2025}:

\begin{equation}
    \xipar = \frac{4}{30000} \approx 1{,}333 \times 10^{-4}
    \label{eq:xi_param}
\end{equation}

Dieser Parameter definiert das Verhältnis zwischen Planck-Skala und kosmischer Skala.
Es gibt \textbf{keine freien Parameter} -- der Grundzustand der Theorie ist durch 
$\xipar$ vollständig bestimmt.

\subsection{T0-Qubits}

In der T0-Theorie emergieren Qubits natürlich aus der Zeit-Masse-Dualität 
\cite{pascher_quantum_2025}:

\begin{equation}
    |\text{T0-Qubit}\rangle = \alpha \, |T{\uparrow}, m{\downarrow}\rangle 
    + \beta \, |T{\downarrow}, m{\uparrow}\rangle
    \label{eq:t0_qubit}
\end{equation}

mit der Bindungsbedingung $T \cdot m = 1$. Die Spinzustände sind nicht willkürlich 
gewählt -- sie sind durch die fundamentale Geometrie erzwungen:

\begin{itemize}
    \item $|T{\uparrow}, m{\downarrow}\rangle$: Hohe Zeit / niedrige Masse 
    \item $|T{\downarrow}, m{\uparrow}\rangle$: Niedrige Zeit / hohe Masse
\end{itemize}

\subsection{Die fraktale Korrektur}

Die fraktale Korrekturgröße $\Kfrak$ spielt in T0 eine zentrale Rolle: 
$\Kfrak^{3/2}$ erklärt die $\sim 2\%$-Abweichung beim anomalen magnetischen 
Moment des Myons (g-2-Anomalie) \cite{pascher_g2_2025}. Dies ist strukturell 
analog zum Annealing-Mechanismus in der CIM.

\clearpage

\section{Sechs Strukturelle Parallelen: T0 und CIM}
\label{sec:parallels}

\subsection{Parallele 1: Determinismus hinter scheinbarer Stochastik}

\begin{table}[h]
    \centering
    \caption{Parallele 1: Determinismus}
    \label{tab:parallel1}
    \resizebox{\textwidth}{!}{\begin{tabular}{@{}lll@{}}
        \toprule
        Aspekt & Coherent Ising Machine & T0-Theorie \\
        \midrule
        Oberfläche & Scheinbar stochastische Suche & Scheinbar probabilistische QM \\
        Tiefenstruktur & Deterministisch (OPO-Feldgl.) & Deterministisch ($\mfield$-Geometrie) \\
        Quelle der ``Zufälligkeit'' & Unbekannte Anfangsbedingungen & Unbekannte Feldkonfiguration \\
        Konsequenz & Lösung ist physikalisch erzwungen & Messergebnis geometrisch bestimmt \\
        \bottomrule
    \end{tabular}}
\end{table}

Beide Systeme ersetzen fundamentalen Probabilismus durch versteckten Determinismus:
In der CIM wird die Lösung nicht berechnet, sie \emph{entsteht}.
In T0 wird das Messergebnis nicht zufällig gewählt, es ist durch das Massenfeld 
$\mfield$ bereits bestimmt.

\begin{keyresult}[Parallele 1]
    Scheinbarer Probabilismus $\equiv$ versteckter Determinismus durch geometrische 
    Feldstruktur -- in CIM \emph{und} T0.
\end{keyresult}

\subsection{Parallele 2: Grundzustand als geometrisches Minimum}

Die CIM löst Optimierung durch physikalische Relaxation:
\begin{equation}
    H_{\text{Ising}} \rightarrow \text{Minimum} \equiv \text{Optimierungslösung}
    \label{eq:cim_groundstate}
\end{equation}

T0 beschreibt das Universum als bereits im Grundzustand seiner eigenen Geometrie:
\begin{equation}
    \xipar = \frac{4}{30000} \equiv \text{geometrisches Minimum} \Rightarrow 
    \text{alle Naturkonstanten}
    \label{eq:t0_groundstate}
\end{equation}

In beiden Fällen ist die ``Lösung'' keine Berechnung -- sie ist der physikalisch 
natürliche Ruhezustand des Systems. \textbf{Optimierung ist Relaxation, nicht Iteration.}

\subsection{Parallele 3: Qubit-Struktur}

\begin{table}[h]
    \centering
    \caption{Parallele 3: Vergleich der Qubit-Strukturen}
    \label{tab:parallel3}
    \resizebox{\textwidth}{!}{\begin{tabular}{@{}lll@{}}
        \toprule
        Aspekt & CIM-Spin & T0-Qubit \\
        \midrule
        Träger & Phase des Laserpulses & Zeit-Masse-Relation \\
        Spin $+1$ & Phase $0°$ & $|T{\uparrow}, m{\downarrow}\rangle$ \\
        Spin $-1$ & Phase $180°$ & $|T{\downarrow}, m{\uparrow}\rangle$ \\
        Constraint & Extern (Laserpumpe) & Intrinsisch: $T \cdot m = 1$ \\
        Superposition & Kohärente Phasenüberlagerung & $\alpha|T{\uparrow},m{\downarrow}\rangle 
        + \beta|T{\downarrow},m{\uparrow}\rangle$ \\
        Kollaps & Phasenmessung & Geometrische Feldprojektion \\
        \bottomrule
    \end{tabular}}
\end{table}

\begin{important}[Fundamentalunterschied]
    Der T0-Spin ist \emph{fundamentaler} als der CIM-Spin: Er trägt seine eigene 
    Constraint ($T \cdot m = 1$) intrinsisch in sich. Ein CIM-Spin muss extern auf 
    $\pm 1$ begrenzt werden; ein T0-Spin ist es durch die Geometrie der Raumzeit selbst.
\end{important}

\subsection{Parallele 4: Die Kopplungsmatrix $J_{ij}$}
\label{subsec:coupling}

Dies ist die tiefste und physikalisch bedeutsamste Parallele:

\begin{table}[h]
    \centering
    \caption{Parallele 4: Kopplungsmatrix -- der entscheidende Unterschied}
    \label{tab:parallel4}
    \resizebox{\textwidth}{!}{\begin{tabular}{@{}lll@{}}
        \toprule
        Aspekt & Klassische CIM & T0-Ising Machine (Konzept) \\
        \midrule
        Kopplungen $J_{ij}$ & Extern programmiert & Geometrisch emergiert \\
        Freiheitsgrade & $J_{ij}$ frei wählbar & Eindeutig durch Feldgeometrie \\
        Physikalische Bedeutung & Mathematisches Hilfsmittel & Reale Wechselwirkungen \\
        Universalität & Beliebige Probleme kodierbar & Physikalisch sinnvolle Probleme \\
        Vorteil & Maximale Flexibilität & Keine Kodierungsfehler möglich \\
        \bottomrule
    \end{tabular}}
\end{table}

In einer T0-basierten Ising Machine gilt $J_{ij}$ \emph{nicht} als freier Parameter, 
sondern emergiert direkt aus dem intrinsischen Zeitfeld:

\begin{equation}
    J_{ij}^{(\text{T0})} = \int \Tfield_i \cdot \Tfield_j \cdot G(x_i, x_j) \, d^3x
    \label{eq:t0_coupling}
\end{equation}

wobei $G(x_i, x_j)$ der geometrische Propagator des T0-Massenfeldes ist.

\subsection{Parallele 5: Fraktale Korrektur als Annealing-Mechanismus}

\begin{table}[h]
    \centering
    \caption{Parallele 5: Fraktale Korrektur und Annealing}
    \label{tab:parallel5}
    \resizebox{\textwidth}{!}{\begin{tabular}{@{}lll@{}}
        \toprule
        Aspekt & Simulated Annealing (CIM) & T0 Fraktale Korrektur \\
        \midrule
        Parameter & Temperatur $T_{\text{ann}}$ & Korrekturgröße $\Kfrak$ \\
        Funktion & Verhindert lokale Minima & Erklärt systematische Abweichungen \\
        Wirkung & Abkühlen $\rightarrow$ Grundzustand & Skalierung $\rightarrow$ Präzision \\
        Beispiel & MaxCut-Optimierung & g-2-Anomalie: $\Kfrak^{3/2}$ erklärt $\sim 2\%$ \\
        Mechanismus & Thermische Fluktuationen & Fraktale Selbstähnlichkeit \\
        \bottomrule
    \end{tabular}}
\end{table}

Die fraktale Korrekturgröße $\Kfrak^{3/2}$ in T0 spielt die analoge Rolle des 
Annealing-Temperaturparameters: sie verhindert, dass das System in einem lokalen 
Minimum der geometrischen Energie gefangen bleibt.

\subsection{Parallele 6: Toroidale Topologie als optimale CIM-Geometrie}

T0 beschreibt das Universum als 4D-Torus-Kristall \cite{pascher_torus_2025}. 
Diese Topologie hat direkte Implikationen für eine T0-basierte Ising Machine:

\begin{itemize}
    \item \textbf{Toroidale Randbedingungen}: Kein Randproblem bei der Optimierung -- 
    alle Spins haben gleich viele Nachbarn
    \item \textbf{Periodische Struktur}: Natürliche Implementierung 
    translationsinvarianter Kopplungen
    \item \textbf{4D statt 3D}: Reichere Kopplungsstruktur als klassische 2D/3D 
    CIM-Arrays
    \item \textbf{Kristallsymmetrie}: Diskrete Symmetriegruppen $\rightarrow$ 
    natürliche Spin-Gitter ohne externe Strukturvorgabe
\end{itemize}

\begin{keyresult}[Parallele 6]
    Eine T0-Ising Machine würde auf der \emph{natürlichen Geometrie des Universums} 
    operieren -- nicht auf einem künstlich konstruierten Graphen.
\end{keyresult}

\clearpage

\section{Die T0-Ising Machine: Konzeptuelle Formulierung}
\label{sec:t0machine}

\subsection{Gesamtstruktur}

Auf Basis der sechs identifizierten Parallelen lässt sich eine T0-basierte 
Ising Machine konzeptuell definieren:

\begin{table}[h]
    \centering
    \caption{Gesamtvergleich: Klassische CIM vs.\ T0-Ising Machine}
    \label{tab:t0machine}
    \resizebox{\textwidth}{!}{\begin{tabular}{@{}lll@{}}
        \toprule
        Komponente & Klassische CIM & T0-Ising Machine \\
        \midrule
        Spin-Träger & DOPO-Laserpulse (Phase) & Zeitfeld-Zustände $\Tfield$ \\
        Spin-Werte & $\pm 1$ (Phase $0°/180°$) & $|T{\uparrow},m{\downarrow}\rangle$, 
        $|T{\downarrow},m{\uparrow}\rangle$ \\
        Kopplungen & Extern programmiert & Geometrisch aus $\mfield$ emergiert \\
        Topologie & Beliebiger Graph & 4D-Torus-Kristall \\
        Grundzustand & Gesuchte Lösung & Geometrischer Fixpunkt $\xipar$ \\
        Annealing & Extern (Temperatur) & Intern ($\Kfrak$-Korrektur) \\
        Determinismus & Versteckt deterministisch & Explizit deterministisch \\
        \bottomrule
    \end{tabular}}
\end{table}

\subsection{Formale Korrespondenz}

Die vollständige Abbildung zwischen T0-Geometrie und Ising-Formalismus lautet:

\begin{align}
    \sigma_i &\leftrightarrow \operatorname{sign}\!\left(\Tfield_i - \bar{T}\right)
    &&\text{[Spin aus Zeitfeld]}
    \label{eq:t0_spin_map}\\
    J_{ij} &\leftrightarrow \langle \Tfield_i \cdot \Tfield_j \rangle_{\xipar}
    &&\text{[Kopplung aus $\xipar$-Geometrie]}
    \label{eq:t0_j_map}\\
    h_i &\leftrightarrow \frac{\partial \xipar}{\partial \Tfield_i}
    &&\text{[Ext. Feld aus $\xipar$-Gradient]}
    \label{eq:t0_h_map}\\
    H_{\text{Ising}} &\leftrightarrow E_{\text{T0}} = \int |\nabla \Tfield|^2 \, d^4x
    &&\text{[Energiefunktionale identisch]}
    \label{eq:t0_energy_map}
\end{align}

Das Minimum des T0-Energiefunktionals \eqref{eq:t0_energy_map} ist identisch mit 
dem Grundzustand des assoziierten Ising-Hamiltonians -- mit dem entscheidenden 
Unterschied, dass in T0 alle Größen ($J_{ij}$, $h_i$, $H$) aus einem einzigen 
Parameter $\xipar$ emergieren.

\subsection{Die zentrale Konsequenz}

\begin{revolutionary}[Fundamentale Implikation der T0-Ising Machine]
    Wenn T0 stimmt, dann ist jedes physikalisch realisierbares Optimierungsproblem 
    bereits als geometrische Struktur im Massenfeld $\mfield$ kodiert. Eine 
    T0-Ising Machine würde nicht optimieren -- sie würde eine physikalische Realität 
    \emph{lesen}. Die Lösung ist nicht das Ergebnis einer Suche, sondern die 
    Enthüllung einer bereits bestehenden geometrischen Wahrheit.
\end{revolutionary}

\clearpage

\section{Offene Fragen und Forschungsrichtungen}
\label{sec:questions}

\subsection{Theoretische Offenheit}

\begin{itemize}
    \item Wie lautet der explizite geometrische Propagator $G(x_i, x_j)$ 
    im T0-Massenfeld? Welche Symmetrien erzwingt die Torus-Topologie?
    
    \item Gibt es eine natürliche Beschränkung, welche NP-harten Probleme 
    durch T0-Geometrie repräsentierbar sind -- oder sind alle Probleme 
    prinzipiell kodierbar?
    
    \item Wie verhält sich $\Kfrak$ als Annealing-Parameter bei verschiedenen 
    Problemklassen? Gibt es universelle Skalierungsgesetze?
    
    \item Welche Verbindung besteht zwischen der Torus-Struktur des T0-Universums 
    und optimalen CIM-Topologien (toroidale vs.\ Möbiusband-Kopplungen)?
\end{itemize}

\subsection{Experimentelle Richtungen}

\begin{itemize}
    \item \textbf{Kopplungsstruktur-Vergleich}: Vergleich der Kopplungsmatrix 
    bestehender CIMs mit T0-Feldpropagator-Vorhersagen. Gibt es systematische 
    Abweichungen, die auf geometrische Einschränkungen hinweisen?
    
    \item \textbf{Toroidale CIM-Topologien}: Zeigen toroidale CIM-Architekturen 
    gegenüber offenen Graphen Leistungsvorteile, die über kombinatorische 
    Erklärungen hinausgehen?
    
    \item \textbf{T0-Qubit-Test}: Zeigen die in \eqref{eq:t0_qubit} definierten 
    T0-Qubits charakteristische Signaturen an bestehenden photonischen 
    Quantencomputing-Experimenten?
    
    \item \textbf{Testbare Vorhersage}: In einem CIM mit $N$ Spins sollte der 
    T0-Korrekturfaktor $\mathcal{D}(N) = \exp(-\xipar \ln N / \Dfrak)$ ein 
    messbares $\ln(N)$-skalierendes Residuum erzeugen, das sich von 
    Standard-Fehlermodellen ($1/N$ oder $e^{-N}$) qualitativ unterscheidet.
\end{itemize}

\subsection{Verbindung zu Scramblon-Physik}

Das vorliegende Dokument schließt an \cite{pascher_scramblons_2026} an, das die 
T0-Verbindung zur Scramblon-Physik analysiert. Die Kombination beider Ansätze 
ergibt ein vollständiges Bild:

\begin{itemize}
    \item \textbf{Scramblon-Theorie}: Dynamische Beschreibung -- wie Quanteninformation 
    in einem CIM-ähnlichen System scrambled
    \item \textbf{T0-Ising Machine}: Geometrische Beschreibung -- welche fundamentalen 
    Kopplungen aus der Raumzeitstruktur folgen
    \item \textbf{Synthese}: Scramblons beschreiben die \emph{Dynamik} der 
    Informationsausbreitung; T0 beschreibt die \emph{Geometrie} der Kopplungsstruktur
\end{itemize}

\clearpage

\section{Zusammenfassung}
\label{sec:summary}

Wir haben sechs strukturelle Parallelen zwischen Coherent Ising Machines und der 
T0-Theorie identifiziert und systematisch analysiert:

\begin{table}[h]
    \centering
    \caption{Zusammenfassung: Sechs Parallelen T0 $\leftrightarrow$ CIM}
    \label{tab:summary}
    \resizebox{\textwidth}{!}{\begin{tabular}{@{}clll@{}}
        \toprule
        \# & Parallele & CIM & T0 \\
        \midrule
        1 & Determinismus & Versteckt (OPO) & Explizit (Geometrie) \\
        2 & Grundzustand & Gesuchte Lösung & Fixpunkt $\xipar$ \\
        3 & Qubit-Struktur & Phase $0°/180°$ & $T \cdot m = 1$-Zustände \\
        4 & Kopplungen $J_{ij}$ & Extern programmiert & Geometrisch emergiert \\
        5 & Annealing & Externe Temperatur & $\Kfrak$-Korrektur \\
        6 & Topologie & Beliebiger Graph & 4D-Torus \\
        \bottomrule
    \end{tabular}}
\end{table}

T0 liefert nicht nur eine konzeptuelle Parallele zur CIM -- es bietet eine 
fundamentalere Begründung: Optimierung ist keine Berechnung, die auf einem 
künstlich konstruierten System ausgeführt wird, sondern physikalische Relaxation 
in eine geometrische Wahrheit, die bereits im Massenfeld des Universums 
eingeschrieben ist.

\begin{keyresult}[Zentrales Ergebnis]
    Die formale Korrespondenz
    \begin{equation}
        H_{\text{Ising}} \leftrightarrow E_{\text{T0}} = \int |\nabla \Tfield|^2 \, d^4x
    \end{equation}
    zeigt: Ising-Optimierung und T0-Energieminimierung sind mathematisch äquivalente 
    Formulierungen desselben geometrischen Prinzips -- mit dem entscheidenden Unterschied, 
    dass in T0 \emph{alle} Kopplungen parameterlos aus $\xipar = 4/30000$ emergieren.
\end{keyresult}

\begin{thebibliography}{99}

\bibitem{inagaki_cim_2016}
T.~Inagaki et al.,
\newblock A coherent Ising machine for 2000-node optimization problems,
\newblock \emph{Science} \textbf{354}, 603 (2016).

\bibitem{mcmahon_cim_2016}
P.~L.~McMahon et al.,
\newblock A fully programmable 100-spin coherent Ising machine with all-to-all connections,
\newblock \emph{Science} \textbf{354}, 614 (2016).

\bibitem{yamamoto_cim_2017}
Y.~Yamamoto et al.,
\newblock Coherent Ising machines -- optical neural networks operating at the quantum limit,
\newblock \emph{npj Quantum Inf.} \textbf{3}, 49 (2017).

\bibitem{honjo_100k_cim_2021}
T.~Honjo et al.,
\newblock 100,000-spin coherent Ising machine,
\newblock \emph{Sci.\ Adv.} \textbf{7}, eabh0952 (2021).

\bibitem{pascher_t0_2025}
J.~Pascher,
\newblock T0-Theorie: Zeit-Masse-Dualität als fundamentales Prinzip,
\newblock T0-Dokumentationsreihe, Dokument 003 (2025).

\bibitem{pascher_xi_2025}
J.~Pascher,
\newblock Der geometrische Ursprungsparameter $\xi = 4/30000$,
\newblock T0-Dokumentationsreihe, Dokument 009 (2025).

\bibitem{pascher_quantum_2025}
J.~Pascher,
\newblock T0-Quantencomputing: Qubits und Logikgatter,
\newblock T0-Dokumentationsreihe, Dokument 147 (2025).

\bibitem{pascher_g2_2025}
J.~Pascher,
\newblock Fraktale Korrekturen und die g-2-Anomalie,
\newblock T0-Dokumentationsreihe, Dokument 018 (2025).

\bibitem{pascher_torus_2025}
J.~Pascher,
\newblock Das Universum als 4D-Torus-Kristall: Geometrische Kosmologie,
\newblock T0-Dokumentationsreihe, Dokument 026 (2025).

\bibitem{pascher_scramblons_2026}
J.~Pascher,
\newblock T0-Projektion auf Scramblon-Physik,
\newblock T0-Dokumentationsreihe, Dokument 148 (2026).

\end{thebibliography}
